\section{Reach to other worlds}

The results sort the adapted models' parts by how far they moved the Moon, and the classes of repair sort
the same parts by how far they travel. The general repairs, which another world takes unchanged, moved the
global mean by a few kelvin at most and moved where sunlight and rain fall; the calibrations and the
Earth-tuned cloud settings, which another world must remake, moved it most. The general repairs are the
synodic Sun clock, the fixed orbit and whole model days, convective cover per Earth day, cloud water by
pressure, gravity throughout CM1 with the fall-speed scaling, CM1's Sun for any day and its deeper soil, the
lifted breakdown cap and the scaled sub-step, SWAN's knee, and the ice optics. Leakage and the methods that
force a slice of a GCM travel with their inputs, a conductivity column or a GCM's fields, and the leader's
crossing travels as a bound set between two of WRF-ELEC's rules; the curved-air terms travel as a method whose
table is computed again for each column; and the calibrations, the Rayleigh and water multipliers, the water continuum, the longwave
correction and the gas-band fits, hold near the case they were fitted on. Another world therefore takes the
general repairs unchanged and needs its calibrations, its ledger and its settings made again, with tools
that come with the repairs: the reference model, the ledger and offline copies of PlaSim's
radiation.\claim{W01}

\textbf{The defects in other users' runs.} Three of ExoPlaSim's defects act well beyond the Moon. Its Sun
clock, in any non-synchronous run whose rotation differs from one day, skips back $1/(N-1)$ of a turn each
model day, where $N$ counts the planet's rotations in a 360-day year, rounded to a multiple of 12 and floored
at 12: about 1\textdegree\ near Earth's rotation, 2\textdegree\ at a two-day rotation, 10\textdegree\ at ten
days and 33\textdegree\ at any rotation slower than 20 days.\claim{W02} Its convective-cover unit adds
$0.125\ln(\text{day}/1\,\text{d})$ to the cover in any run whose model day differs from one Earth day,
synchronous runs included, since ExoPlaSim's synchronous mode fixes the Sun in the sky and keeps the model
day its rotation period sets: about $+0.23$ for TRAPPIST-1e's rotation of 6.1 days and $+0.31$ for Proxima
Centauri~b's 11.2 days. Its default orbit gives Earth's seasons to any run that sets an obliquity or an
eccentricity without fixing the orbit.\claim{W06}

\textbf{Where Earth values break.} Against the ratios of Figure~\ref{fig:regimes}, each Earth value fails
along its own direction. PlaSim's horizon factor of 35 is right only at Earth's scale height over radius; an
isothermal Chapman factor gives 22 for present Mars, 24.5 for Venus, 14 for Titan and 7.2 for the Moon, so the
slant path errs for each and most for the tallest air, while PlaSim's broadband fits, made near Earth's
column, meet columns ten times Earth's on Titan, a hundred times on Venus and seven times on the Moon.
Slingo's cover, with its unit repaired, remains a fit made on Earth's convection, and PlaSim's cloud
water, written in metres, errs at any gravity other than Earth's, too dry aloft under weaker gravity and too
wet under stronger.\claim{W02} Released CM1 fails at any gravity but Earth's, under any day but Earth's
24 hours, and in any air whose carbon dioxide differs from Earth's value for the calendar year; SWAN, in the
configuration used, loses gravity similarity through one constant. Taken model by model, the adapted CM1
runs storms in Earth-like air at any gravity on its Morrison and NSSL paths, SWAN and SWASH run water seas at
any gravity by similarity with Earth's empirical laws, and ExoPlaSim, bounded by its radiation fits and its
cloud, has the narrowest reach of the three. Table~\ref{tab:reach} sets out, property by property, where the
released models break, what the adaptations now reach and what lies beyond them. Of the found worlds in
Figure~\ref{fig:regimes}, only Earth lies within that reach; the rest of it lies among constructed worlds and
among exoplanets whose air is taken to be close to Earth's.

\textbf{Within reach as adapted.} The breathable Moon under other dimmings runs in all four models with the
repairs as they stand, as long as its columns stay near the fitted ones: the split of sunlight holds to 0.7
points on a 303~K column with 338~kg\,m$^{-2}$ of vapour, while the longwave correction's accuracy near 307~K
is unchecked. At 1.0~atm, a case yet to be run, or behind another shield that blocks the ultraviolet, it
needs its calibrations refitted against the reference and its cloud settings checked; a shield that lets the
ultraviolet through needs ozone, which the CM1 build sets to zero. Slowly rotating and synchronous planets
with Earth's gravity and air, common subjects of exoplanet GCMs \autocite{yang2014strong,way2018climates},
should need only ExoPlaSim's general repairs, though none has been run: the released radiation stays near the
column it was fitted on as long as the climate stays near Earth's warmth. For rotations slower than 20 days,
ExoPlaSim's model day departs from the solar day unless the year holds about twelve rotations, so output kept
in whole model days drifts against the Sun; analyses that sort output by the Sun's position avoid the
drift.\claim{W03}

\textbf{Within reach with a new column.} A constructed world with breathable air on a small body, such as
Mars held at 1~bar, has a scale height over radius five times Earth's, 2.6 times Earth's air over each square
metre and 0.38 of Earth's gravity. It needs the general repairs, a ledger of its own column and calibrations
refitted to it, and at its 1~bar level it stays inside RRTMG's pressure range, though ground standing below
that level takes RRTMG past its tables, as the lunar seas do. Earth-like air on a larger planet comes within
reach the same way, with the cloud water's error running the other way, and a breathable moon of a giant
planet needs, besides its column, the downward infrared that PlaSim now takes at its top for the host's
glow.\claim{W04}

\textbf{Beyond reach.} Titan's nitrogen--methane air at 94~K, present Mars's thin carbon dioxide and Venus's
92~bar lie outside what PlaSim's and RRTMG's radiation and Morrison's water-only microphysics can represent;
their general repairs would hold, and their radiation, condensation and clouds would have to be
built.\claim{W05} Four limits bound most of the reach of the adapted models: compositions beyond nitrogen,
oxygen, argon, water and modest carbon dioxide; surface pressures past RRTMG's tables, already extrapolated
over the lowest 7.5~km at 1.2~atm; scale heights past a few hundredths of the radius, where thin-shell,
constant-gravity dynamics lose the 5--15\% fall of gravity across the air; and empirical laws of cloud, sea
and charge used far from the Earth conditions they were fitted on. Table~\ref{tab:reach} adds the narrower
ones.

\begin{table}[tbp]
\caption{Reach by world property: where the released ExoPlaSim, CM1, SWAN and SWASH break, what the
adaptations reach, and what lies beyond them.}
\label{tab:reach}
\footnotesize
\setlength{\tabcolsep}{4pt}
\renewcommand{\arraystretch}{1.08}
\begin{tabularx}{\textwidth}{@{}>{\raggedright\arraybackslash}p{0.15\textwidth}LLL@{}}
\toprule
Property & Breaks in the released models & Reached by the adaptations & Beyond them \\
\midrule
Gravity other than Earth's & CM1 throughout; PlaSim's cloud water in metres; SWAN's wind-input knee & Any gravity on CM1's Morrison and NSSL paths; cloud water by pressure; SWAN and SWASH by similarity & CM1's other schemes at 9.81; drop limits and breakup \\
Solar day other than Earth's & ExoPlaSim's Sun clock off a one-day rotation; its convective-cover unit; the output calendar; CM1's 24-hour Sun and diurnal soil; WRF-ELEC lacking clear-air conduction & Synodic clock; cover per Earth day; whole model days; CM1's Sun for any day; soil 1.55~m deep for days from a few Earth days to the Moon's; leakage & CM1's seasons; soil for days far longer than the Moon's; ExoPlaSim's model day for slow rotators; Slingo's fit itself \\
Tilt and orbit & ExoPlaSim's default puts in Earth's orbit & The configured orbit, held fixed & CM1's equinox Sun \\
Air tall against the radius & Flat, sunlit-only columns; PlaSim's horizon factor of 35; Earth's $\sigma$ levels & Curved-air terms and slant factors from a ledger of the world's column & Thin-shell, constant-gravity dynamics; one clear column for all; the vertical grid \\
Air heavier than Earth's & Broadband shortwave and longwave fits; Rayleigh scaled to a 5772~K Sun; RRTMG's tables past 1,053.63~hPa & Multipliers, gas bands and longwave by height, refitted to the reference & RRTMG past 1.05~bar; multi-bar surfaces \\
Composition & RRTMG's calendar CO$_2$, trace gases and ozone; PlaSim's sunlight lacking CO$_2$ and O$_2$ & Nitrogen--oxygen air with O$_2$ and CO$_2$ near the design air's, refitted otherwise; ozone-free in CM1 & CO$_2$- or CH$_4$-dominated air; condensables other than water \\
Cold cloud & Liquid optics for all cloud; saturation over liquid & Ice optics below $-38$\,\textdegree C & Ice saturation; crystal size with gravity; optics below 180~K \\
Storm electricity & The 180~kV\,m$^{-1}$ cap; the $-7$\,\textdegree C ground rule; the 0.75~s sub-step & Density-scaled breakdown; the leader's crossing; a scaled sub-step & Charging in little cloud water; breakdown in other gases \\
Seas & The wind-input knee; the air's density; breaker parameter read uninitialised & Similarity at any gravity; air density per build; breaking kept on & Earth's drag, growth and breaking laws \\
\bottomrule
\end{tabularx}
\end{table}
